\documentclass[12pt,a4paper]{article}

\usepackage[UTF8]{ctex}
\usepackage{amsmath,amssymb,amsthm}
\usepackage{mathtools}
\usepackage{graphicx}
\usepackage{tikz}
\usepackage{hyperref}
\usepackage{geometry}
\usepackage{enumitem}
\usepackage{tcolorbox}
\usepackage{booktabs}
\usepackage{xfrac}
\tcbuselibrary{theorems}

\geometry{margin=2.5cm}

\theoremstyle{plain}
\newtheorem{theorem}{定理}[section]
\newtheorem{proposition}[theorem]{命题}
\newtheorem{corollary}[theorem]{推论}
\newtheorem{lemma}[theorem]{引理}

\theoremstyle{definition}
\newtheorem{definition}[theorem]{定义}
\newtheorem{example}[theorem]{例}
\newtheorem{problem}[theorem]{问题}

\theoremstyle{remark}
\newtheorem{remark}[theorem]{注记}
\newtheorem{conjecture}[theorem]{猜想}

\title{\bfseries 接触图、Tur\'an 定理与平面密铺\\[4pt]
\large ——可密铺全平面的多边形之图论刻画}
\author{}
\date{\today}

\begin{document}

\maketitle

\begin{abstract}
本文从接触图（contact graph）的视角研究平面密铺问题。
将每块瓷砖视为顶点、共享边界视为边，构造平面图。
利用 Euler 公式与图论约束，我们直接从接触图推导出：
（1）边对边的单一凸多边形密铺中，多边形边数 $k \le 6$；
（2）正多边形密铺的顶点类型完全由丢番图方程 $\sum 1/k_i = (m-2)/2$ 决定；
（3）Tur\'an 定理在直接应用中虽不紧，但``平面 Tur\'an 数''——即禁用于结构下平面图的最大边数——为多类型瓷砖的接触模式提供了非平凡的约束。
\end{abstract}

\tableofcontents

% ============================================================
\section{接触图：从铺砌到图论}
% ============================================================

\subsection{定义}

\begin{definition}[接触图]
设 $\mathcal{T} = \{T_1, \dots, T_n\}$ 为平面的一个铺砌（瓷砖内部两两不交，并集覆盖区域）。
$\mathcal{T}$ 的\textbf{接触图} $G = G(\mathcal{T})$ 定义为：
\begin{itemize}
    \item $V(G) = \{v_1, \dots, v_n\}$，$v_i$ 对应瓷砖 $T_i$；
    \item 边 $v_i v_j \in E(G)$ 当且仅当 $\partial T_i \cap \partial T_j$ 包含长度为正的线段。
\end{itemize}
\end{definition}

\begin{proposition}[接触图的基本性质]\label{prop:basic}
对于平面铺砌的接触图 $G$：
\begin{enumerate}[label=(\roman*)]
    \item $G$ 是\textbf{平面图}（可取每块瓷砖内部一点为顶点，沿共享边中点连线为边）；
    \item 若铺砌为边对边（edge-to-edge），且所有瓷砖为 $k$ 边形，则 $G$ 是近似 $k$-正则的（内部顶点度数 $=k$）；
    \item 若全平面密铺，取大块 $n$ 枚瓷砖，边界效应可忽略，视为环面上的图，满足 $V - E + F = 0$。
\end{enumerate}
\end{proposition}

% ============================================================
\section{单一多边形密铺：$k \le 6$ 的推导}\label{sec:monohedral-contact}
% ============================================================

\subsection{Euler 公式推导}

考虑由全等 $k$ 边形构成的边对边平面密铺。取足够大的一块，含 $n$ 枚瓷砖。
将这块粘贴成环面（周期边界条件），则 Euler 公式为：
\begin{equation}\label{eq:euler-torus}
V - E + F = 0,
\end{equation}
其中：
\begin{itemize}
    \item $V = n$ (接触图的顶点数 = 瓷砖数)；
    \item $E = \dfrac{kn}{2}$ (每枚 $k$ 边形有 $k$ 个邻居，每条边被两端共享)；
    \item $F$ = 接触图的面数，对应铺砌中瓷砖相遇的\textbf{顶点}。
\end{itemize}

代入 \eqref{eq:euler-torus}：
\begin{equation}\label{eq:F}
F = E - V = \frac{kn}{2} - n = n\left(\frac{k}{2} - 1\right).
\end{equation}

令 $m$ 为每个铺砌顶点处相遇的瓷砖数（即接触图中每个面的度数）。
所有面的度数之和等于 $2E$：
\begin{equation}\label{eq:face-sum}
\sum_{\text{faces } f} \deg(f) = 2E = kn.
\end{equation}

记平均面度 $\overline{m} = \frac{1}{F}\sum \deg(f)$，则：
\begin{equation}\label{eq:m-avg}
\overline{m} = \frac{kn}{F}
= \frac{kn}{n(k/2 - 1)}
= \frac{2k}{k-2}.
\end{equation}

\begin{theorem}[单一多边形密铺的必要条件]\label{thm:k-bound}
若存在全等 $k$ 边形的边对边平面密铺，则 $k \le 6$。
\end{theorem}

\begin{proof}
任何密铺的顶点处至少有 $3$ 枚瓷砖相遇（否则为边界或退化情形）。故 $\overline{m} \ge 3$。
由 \eqref{eq:m-avg}：
\[
\frac{2k}{k-2} \ge 3 \implies 2k \ge 3k - 6 \implies k \le 6.
\]
\end{proof}

\begin{center}
\begin{tabular}{c|c|c|c}
\toprule
$k$ & $\overline{m} = 2k/(k-2)$ & 可行性 & 已知密铺 \\
\midrule
3 & $6$ & \checkmark & 任意三角形可密铺 \\
4 & $4$ & \checkmark & 任意四边形可密铺 \\
5 & $10/3 \approx 3.33$ & \checkmark & 15 类五边形已知 \\
6 & $3$ & \checkmark & 3 类六边形已知 \\
$\ge 7$ & $<3$ & \textcolor{red}{\textbf{×}} & 不存在 \\
\bottomrule
\end{tabular}
\end{center}

\begin{remark}
表中所列的已知分类为 Reinhardt (1918), Kershner (1968), James (1975), Rice (1977),
Stein (1985), Mann--McLoud--Von Derau (2015--2017) 及 Rao (2017) 等人的工作。
五边形密铺共 15 类，六边形共 3 类。$k \ge 7$ 的凸多边形不存在边对边单一密铺，
这一定理 \ref{thm:k-bound} 提供了简洁的图论证明（但需假设边对边条件）。
\end{remark}

\subsection{非边对边密铺的情况}

若放宽边对边条件（允许瓷砖的部分边与多个邻居接触），接触图的度数可能超过 $k$。
此时 $E$ 增大，$\overline{m}$ 减小。但平面图的 $E \le 3V-6$ 仍适用。
详细分析见第 \ref{sec:multi-edge} 节。

% ============================================================
\section{正多边形密铺与 Archimedean 顶点类型}
% ============================================================

\subsection{顶点类型的丢番图方程}

设在一个密铺顶点处有 $m$ 枚正 $k_1, k_2, \dots, k_m$ 边形相遇。
正 $k$ 边形的内角为 $\frac{k-2}{k}\pi$。所有角度之和应为 $2\pi$：
\[
\sum_{i=1}^{m} \frac{k_i-2}{k_i} = 2
\implies \sum_{i=1}^{m} \frac{1}{k_i} = \frac{m-2}{2}.
\]
这正是第 \ref{sec:monohedral-contact} 节中 $\overline{m} = 2k/(k-2)$ 的``局部''版本。

\subsection{枚举结果}

通过计算程序枚举（$k_i \le 12$）得到所有解：

{\small
\begin{center}
\begin{tabular}{c|c|l}
\toprule
$m$ & 目标 $\sum 1/k_i$ & 顶点类型 $(k_1, \dots, k_m)$ \\
\midrule
3 & $1/2$ & $(3,12,12),\;(4,6,12),\;(4,8,8),\;(5,5,10),\;(6,6,6)$ \\
4 & $1$   & $(3,3,4,12),\;(3,3,6,6),\;(3,4,4,6),\;(4,4,4,4)$ \\
5 & $3/2$ & $(3,3,3,3,6),\;(3,3,3,4,4)$ \\
6 & $2$   & $(3,3,3,3,3,3)$ \\
\bottomrule
\end{tabular}
\end{center}
\par}

这些恰好是\textbf{Archimedean 密铺}（均匀密铺）的所有顶点类型。其中：
\begin{itemize}
    \item $(6,6,6)$ 对应正六边形密铺
    \item $(4,4,4,4)$ 对应正方形网格
    \item $(3,3,3,3,3,3)$ 对应正三角形密铺
    \item $(3,6,3,6)$ 为三六边形密铺（trihexagonal），但由 $(3,3,6,6)$ 经不同顺序排列得到
\end{itemize}

\subsection{从顶点类型到接触图参数}

对于每种 Archimedean 密铺，可计算瓷砖类型比例及接触图的平均度数：

{\small
\begin{center}
\begin{tabular}{c|c|c|c}
\toprule
顶点类型 & 瓷砖比例 & 平均边数 $k_{\text{avg}}$ & 预测 $\overline{m}$ \\
\midrule
$3^6$ & 三角形 100\% & $3$ & $6$ \\
$4^4$ & 正方形 100\% & $4$ & $4$ \\
$6^3$ & 六边形 100\% & $6$ & $3$ \\
$3^3.4^2$ & $\triangle\,66.7\%,\;\square\,33.3\%$ & $3.333$ & $5$ \\
$3^4.6$ & $\triangle\,88.9\%,\;$ 六边形 $11.1\%$ & $3.333$ & $5$ \\
$3.6.3.6$ & $\triangle\,66.7\%,\;$ 六边形 $33.3\%$ & $4$ & $4$ \\
$3.4.6.4$ & $\triangle\,33.3\%,\;\square\,50\%,\;$ 六边形 $16.7\%$ & $4$ & $4$ \\
$4.8^2$ & $\square\,50\%,\;$ 八边形 $50\%$ & $6$ & $3$ \\
\bottomrule
\end{tabular}
\end{center}
\par}

所有类型均满足 $\overline{m} = 2k_{\text{avg}}/(k_{\text{avg}}-2)$，验证了公式 \eqref{eq:m-avg} 的普适性。

% ============================================================
\section{Tur\'an 定理与铺砌的联系}
% ============================================================

\subsection{直接应用的局限性}

Tur\'an 定理：$n$ 个顶点、不含 $K_{r+1}$ 的图边数至多为 $t(n,r) \approx \left(1 - \frac{1}{r}\right)\frac{n^2}{2}$。

密铺接触图是平面图，故不含 $K_5$ 及 $K_{3,3}$，对应 $r=4$。
Tur\'an 界 $t(n,4) \approx \frac{3}{8}n^2$ 是 $O(n^2)$ 量级。
然而平面图的边数仅为 $O(n)$（至多 $3n-6$），远小于 Tur\'an 界。
因此，纯 Tur\'an 定理对密铺的直接约束\textbf{不紧}。

\begin{center}
\begin{tabular}{c|r|r|r}
\toprule
$k$ & 实际边数 $E=kn/2$ ($n=1000$) & Tur\'an 界 $t(1000,4)$ & 平面界 $3n-6$ \\
\midrule
3 & 1,500 & 375,000 & 2,994 \\
4 & 2,000 & 375,000 & 2,994 \\
5 & 2,500 & 375,000 & 2,994 \\
6 & 3,000 & 375,000 & 2,994 \\
\bottomrule
\end{tabular}
\end{center}

\subsection{真正有意义的联系：平面 Tur\'an 数}

Tur\'an 定理的精髓在于：\textbf{禁止某个子结构 → 控制整体密度 → 极值图具有部结构}。
当我们将此思想\textbf{限制在平面图类中}时，得到的是\textbf{平面 Tur\'an 数}（planar Tur\'an number）：

\begin{definition}[平面 Tur\'an 数]
$\mathrm{ex}_{\mathcal{P}}(n, H) = \max\{|E(G)| : G \text{ 是 } n \text{ 个顶点的平面图，且不含 } H \text{ 作为子图}\}$。
\end{definition}

这与密铺的联系非常自然：密铺的接触图是平面图；若进一步禁止某种接触模式（如``同类瓷砖不相邻''），
则接触图属于特定的平面 Tur\'an 极值问题。

\subsection{同类瓷砖不相邻：$r$ 部平面图}

假设密铺中有 $r$ 种类型的瓷砖，且\textbf{同类瓷砖彼此不相邻}。
则接触图为 $r$ 部平面图（各部分为独立集）。

\begin{proposition}[$r$ 部平面图的最大边数]
记 $f_r(n)$ 为 $n$ 顶点的 $r$ 部平面图最大边数：
\begin{itemize}
    \item $r=2$（二部平面图）：$f_2(n) \le 2n-4$（对 $n \ge 3$）；
    \item $r=3$（三部平面图）：$f_3(n)$ 已知为某线性函数，且受限于不含 $K_{3,3}$；
    \item $r \ge 4$：$f_r(n) \le 3n-6$（等同于一般平面图界）。
\end{itemize}
\end{proposition}

\begin{remark}
与 Tur\'an 图 $T(n,r)$ 相比：$T(n,2) = K_{\lfloor n/2\rfloor, \lceil n/2\rceil}$ 有 $\approx n^2/4$ 条边，
但二部平面图至多 $2n-4$ 条边——差了整整一个量级！
这正是\textbf{平面性比 Tur\'an 条件强 $O(n)$ 倍}的体现。
\end{remark}

\begin{example}[三六边形密铺]
在三六边形密铺 (3.6.3.6) 中：三角形只接触六边形，六边形只接触三角形。
接触图是\textbf{二部图}。每枚三角形接触 $3$ 个六边形，每枚六边形接触 $6$ 个三角形。
若总共有 $a$ 个三角形和 $b$ 个六边形，则边数 $E = 3a = 6b$，故 $a = 2b$。
接触图是二部平面图，边数 $E = 3a = 3 \cdot 2b = 6b$。
顶点数 $n = a + b = 3b$。平面二部图最大边数为 $2n-4 = 6b-4$，
实际边数 $6b$（忽略边界效应）恰好接近此上界。
\end{example}

% ============================================================
\section{``重边''与多重接触}\label{sec:multi-edge}
% ============================================================

\subsection{凸多边形密铺中不存在真正的重边}

对于凸多边形的密铺，两枚瓷砖的边界交集 $\partial T_i \cap \partial T_j$ 是两个凸集的交，
依然为凸集，故只能是空集、一个点、或一条线段。不可能出现两枚瓷砖在两个不连通的线段上接触。
因此接触图在简单图意义下\textbf{不存在重边}。

\subsection{推广一：带权接触图}

若不仅关心``是否接触''，还关心接触长度，定义\textbf{带权接触图}：
\[
w(v_i v_j) = \mathcal{H}^1(\partial T_i \cap \partial T_j) = \text{共享边界的长度}.
\]

此时 Tur\'an 型问题变为：\textbf{在禁止某种带权团结构的条件下，总权重的极值}。

\begin{problem}[带权 Tur\'an--密铺问题]
给定 $n$ 枚面积固定的正方形。将其排列使总接触长度最大化，问最大值及极值排列。
\end{problem}

这类似于等周问题的离散版本：固定总面积，最大化内部接触。

\subsection{推广二：非边对边密铺中的多重接触}

在非边对边密铺中，一枚瓷砖可以沿其一条边与多个邻居接触。
此情况下，每枚瓷砖的度数可能超过 $k$（$k$ 边形可以有超过 $k$ 个邻居）。
接触图可能具有更高的平均度数，使 $k \le 6$ 的论证需要修正。

若进一步推广至三维：多面体可以通过多个面接触，形成真正的多重边（multi-edges），
此时图论框架自然地容纳重边图或超图。

% ============================================================
\section{总结与开放问题}
% ============================================================

\subsection{主要结果}

\begin{enumerate}[label=(\arabic*)]
    \item 接触图 + Euler 公式 $\Rightarrow$ 边对边单一凸 $k$ 边形密铺的必要条件 $k \le 6$。
    这是一个纯图论的推导，不需要任何几何论证（除凸性和平面性外）。
    
    \item 正多边形密铺的顶点类型完全由 $\sum 1/k_i = (m-2)/2$ 决定。
    这等价于内角之和 $= 360^\circ$，但接触图框架揭示了它与平均面度公式 $\overline{m} = 2k/(k-2)$ 的统一性。
    
    \item Tur\'an 定理对接触图的直接约束不紧（$O(n^2)$ vs $O(n)$），
    但``平面 Tur\'an 数''——限制在平面图类中的极值——为多类型瓷砖的接触模式提供非平凡约束。
    尤其是同类瓷砖不相邻时，$r$ 部平面图的最大边数远小于一般 $r$ 部图的 Tur\'an 界。
\end{enumerate}

\subsection{开放问题}

\begin{enumerate}[label=\textbf{Q\arabic*}.]

    \item \textbf{平面 Tur\'an 数的几何实现}：
    对于给定的禁用于图 $H$，$\mathrm{ex}_{\mathcal{P}}(n, H)$ 的极值平面图
    能否实现为\textbf{某种密铺的接触图}？哪些平面 Tur\'an 极值图有几何对应？
    
    \item \textbf{五边形密铺的图论分类}：
    对于 $k=5$，$\overline{m} = 10/3$，意味顶点类型为 $3$ 或 $4$ 的组合。
    能否从接触图的角度理解为何恰好有 $15$ 类五边形密铺？
    
    \item \textbf{非边对边密铺的接触图}：
    如何修正 Euler 公式推导以涵盖非边对边情形？此时接触图的平均度数不再等于 $k$。
    
    \item \textbf{带权极值问题}：
    固定总面积的 $n$ 枚多边形，最大化接触总长度。极值排列与等周不等式有何联系？
    
    \item \textbf{三维推广}：
    多面体堆砌的接触图（面接触）满足何种 Euler 型约束？重边在三维中有非平凡实例。
    
    \item \textbf{非周期密铺与无限图}：
    对于非周期（如 Penrose 密铺）的接触图（无限平面图），
    Tur\'an 型密度参数是否有意义？

\end{enumerate}

\vspace{12pt}
\begin{tcolorbox}[title=核心洞察]
密铺问题的图论本质：\textbf{接触图是介于组合与几何之间的桥梁}。
它既承载了平面性这一几何信息，又将``谁与谁相邻''提升为独立的组合对象。
Tur\'an 定理的思想——通过禁止子结构来控制密度——在限制为平面图类后，
为分析瓷砖间的接触模式（尤其是多类型密铺中同类相避的情形）提供了有力的框架。
\end{tcolorbox}

\end{document}
